%======== 'Results' Declarations =========
% Arguments to this command are:
% 1. Environment ID (e.g., theorem)
% 2. Environment name (e.g., Theorem)
% 3. Colour
% 4. Counter increment (e.g., chapter, section)
\newcommand{\makenewresultstyle}[4]{
    \newcounter{#1}[#4]
    \newcommand{\internalcounter}{\csname the#1\endcsname}
    \newenvironment{#1}[1][]
    {
        \stepcounter{#1}
        \ifstrempty{##1}
        {
            \mdfsetup{
                frametitle={
                    \tikz[baseline=(current bounding box.east),outer sep=0pt]
                    \node[anchor=east,rectangle,fill=#3]{\strut #2~\internalcounter};
                }
            }
        }
        {
            \mdfsetup{
                frametitle={
                    \tikz[baseline=(current bounding box.east),outer sep=0pt]
                    \node[anchor=east,rectangle,fill=#3]{\strut #2~\internalcounter:~##1};
                }
            }
    }
    \mdfsetup{
            innertopmargin=10pt,linecolor=#3,
            linewidth=2pt,topline=true,
            frametitleaboveskip=\dimexpr-\ht\strutbox\relax
        }
        \begin{mdframed}[]\relax
    }
    {\end{mdframed}}
}

% \newcommand{\makenewresultstyle}[3]{
%     \declaretheoremstyle[
%         headfont=\normalfont\bfseries,
%         bodyfont=\normalfont,
%         notefont=\normalfont\bfseries\boldmath\itshape,
%         spaceabove=0pt,  % Space between previous paragraph and current block
%         spacebelow=0pt,  % Space between current block and next paragraph
%         mdframed={
%             linewidth=0.75pt,
%             linecolor=#3,
%             skipabove=2.8pt,  % Space between top of block and beginning of coloured frame
%             skipbelow=2.8pt,  % Space between bottom of block and beginning of coloured frame
%             backgroundcolor=#2,
%             usetwoside=false,  % Needed for `leftmargin` and `rightmargin` to work
%             leftmargin=-5pt,
%             rightmargin=-5pt,
%             innerleftmargin=5pt,
%             innerrightmargin=5pt
%         }
%     ]{#1-style}
% }

% \makenewresultstyle{theorem}{DarkSeaGreen2}{DarkSeaGreen4}
% \declaretheorem[name=Theorem,style=theorem-style,within=section]{theorem}
% \renewcommand*{\thetheorem}{\arabic{chapter}.\arabic{section}.\arabic{theorem}}

% \makenewresultstyle{lemma}{AntiqueWhite1}{Bisque4}
% \declaretheorem[style=lemma-style,sibling=theorem]{lemma}

% \makenewresultstyle{proposition}{Honeydew1}{SpringGreen4}
% \declaretheorem[style=proposition-style,sibling=theorem]{proposition}

% \makenewresultstyle{corollary}{Cornsilk1}{Ivory4}
% \declaretheorem[style=corollary-style,sibling=theorem]{corollary}

% \makenewresultstyle{definition}{LightCyan1}{Cyan4}
% \declaretheorem[style=definition-style,sibling=theorem]{definition}

% \makenewresultstyle{axiom}{Thistle2}{Plum4}
% \declaretheorem[style=axiom-style,sibling=theorem]{axiom}

% \newcounter{definition}[section]

\makenewresultstyle{definition}{Definition}{LightCyan1}{section}
\renewcommand*{\thedefinition}{\arabic{chapter}.\arabic{section}.\arabic{definition}}

% \newcounter{definition}[section]
% \renewcommand*{\thedefinition}{\arabic{chapter}.\arabic{section}.\arabic{definition}}
% \newenvironment{definition}[1][]
% {
%     \stepcounter{definition}
%     \ifstrempty{#1}
%     {
%         \mdfsetup{
%             frametitle={
%                 \tikz[baseline=(current bounding box.east),outer sep=0pt]
%                 \node[anchor=east,rectangle,fill=blue!20]{\strut Theorem~\thedefinition};
%             }
%         }
%     }
%     {
%         \mdfsetup{
%             frametitle={
%                 \tikz[baseline=(current bounding box.east),outer sep=0pt]
%                 \node[anchor=east,rectangle,fill=blue!20]{\strut Theorem~\thedefinition:~#1};
%             }
%         }
%    }
%    \mdfsetup{
%         innertopmargin=10pt,linecolor=blue!20,
%         linewidth=2pt,topline=true,
%         frametitleaboveskip=\dimexpr-\ht\strutbox\relax
%     }
%     \begin{mdframed}[]\relax
% }
% {\end{mdframed}}

%======= 'Questions' Declarations ========
% \declaretheoremstyle[
%     headfont=\normalfont\bfseries,
%     bodyfont=\normalfont,
%     spaceabove=5pt,  % Space between previous paragraph and current block
%     prefoothook={\vspace{5pt}},
%     mdframed={
%         linewidth=0.75pt,
%         skipabove=2.8pt,   % Space between top of block and beginning of frame
%         skipbelow=2.8pt,   % Space between bottom of block and beginning of frame
%         usetwoside=false,  % Needed for `leftmargin` and `rightmargin` to work
%         leftmargin=-5pt,
%         rightmargin=-5pt,
%         innerleftmargin=5pt,
%         innerrightmargin=5pt
%     }
% ]{exercise-style}
% \declaretheorem[style=exercise-style,within=chapter]{exercise}
% \renewcommand*{\theexercise}{\arabic{chapter}.\arabic{exercise}}

% \declaretheoremstyle[
%     headfont=\normalfont\bfseries,
%     bodyfont=\normalfont,
%     notefont=\normalfont\bfseries\itshape,
%     spaceabove=10pt,  % Space between previous paragraph and current block
%     spacebelow=10pt,  % Space between current block and next paragraph
% ]{problem-style}
% \declaretheorem[name=Problem,style=problem-style,within=chapter]{problem}
% \renewcommand*{\theproblem}{\arabic{chapter}.\arabic{problem}}

\newcounter{theo}[section]
\newenvironment{theo}[1][]
{
    \stepcounter{theo}
    \ifstrempty{#1}
    {
        \mdfsetup{
            frametitle={
                \tikz[baseline=(current bounding box.east),outer sep=0pt]
                \node[anchor=east,rectangle,fill=blue!20]{\strut Theorem~\thetheo};
            }
        }
    }
    {
        \mdfsetup{
            frametitle={
                \tikz[baseline=(current bounding box.east),outer sep=0pt]
                \node[anchor=east,rectangle,fill=blue!20]{\strut Theorem~\thetheo:~#1};
            }
        }
   }
   \mdfsetup{
        innertopmargin=10pt,linecolor=blue!20,
        linewidth=2pt,topline=true,
        frametitleaboveskip=\dimexpr-\ht\strutbox\relax
    }
    \begin{mdframed}[]\relax
}
{\end{mdframed}}

%====== Miscellaneous Declarations =======
% \declaretheoremstyle[
%     headfont=\normalfont\bfseries,
%     bodyfont=\normalfont,
%     spaceabove=10pt,  % Space between previous paragraph and current block
%     spacebelow=10pt,  % Space between current block and next paragraph
% ]{general-style}
% \declaretheorem[style=general-style,sibling=theorem]{example}
% \declaretheorem[style=general-style,numbered=no]{remark}
